\documentclass[11pt,a4paper]{article}
\usepackage[margin=24mm,headheight=14pt]{geometry}
\usepackage{amsmath,amssymb,amsthm,mathtools,booktabs,array}
\usepackage{lmodern,microtype,fancyhdr}
\usepackage{fontspec,fancyvrb}
\usepackage[hidelinks]{hyperref}
\newfontfamily\leanappendixfont{DejaVuSansMono.ttf}
\newtheorem{theorem}{Theorem}
\newtheorem{lemma}[theorem]{Lemma}
\newtheorem{proposition}[theorem]{Proposition}
\newcommand{\R}{\mathbb R}
\newcommand{\nn}{\R_{\ge0}}
\newcommand{\tr}{\operatorname{tr}}
\newcommand{\diag}{\operatorname{diag}}
\newcommand{\gap}{\operatorname{gap}}
\newcommand{\rev}{\operatorname{rev}}
\newcommand{\code}[1]{\texttt{#1}}
\DeclareMathOperator{\im}{im}
\pagestyle{fancy}
\fancyhf{}
\fancyhead[L]{\small Affine interpretations in two coordinates}
\fancyhead[R]{\small Lucas Valbuena}
\fancyfoot[C]{\thepage}
\renewcommand{\headrulewidth}{0.3pt}
\setlength{\parindent}{1.1em}
\setlength{\parskip}{3pt}
\setlength{\abovedisplayskip}{7pt plus 2pt minus 2pt}
\setlength{\belowdisplayskip}{7pt plus 2pt minus 2pt}
\hypersetup{pdftitle={A two-coordinate obstruction for affine interpretations of the Yolcu--Aaronson--Heule Collatz system},pdfauthor={Lucas Valbuena},pdfsubject={Lean-checked obstruction for strictly monotone affine interpretations}}

\begin{document}
\thispagestyle{plain}
\begin{center}
{\LARGE\bfseries A two-coordinate obstruction for affine\\[3pt]
interpretations of the Yolcu--Aaronson--Heule\\[3pt]
Collatz system\par}
\vspace{8pt}
{\large Lucas Valbuena\par}
\vspace{4pt}
{\normalsize 27 September 2026}
\end{center}
\vspace{2pt}
\begin{abstract}
We prove a restriction on matrix-interpretation proofs for the
eleven-rule Collatz system constructed by Yolcu, Aaronson, and Heule.
For nonnegative real affine maps in two coordinates, assume that every
symbol matrix has first diagonal entry at least one. If all eleven rules
are weakly oriented coefficientwise, every first-coordinate offset gap
is zero. Thus no rule can be removed in a first step using this strictly
monotone interpretation format, regardless of coefficient size. The
same conclusion holds after word reversal and allows arbitrary
admissible boundary maps. The proof includes singular matrices and zero
secondary entries. Lean 4 checks the complete theorem and its supporting
algebra. The relaxed TOP setting and arctic interpretations are outside
the theorem's scope.
\end{abstract}

\section{Introduction}
Yolcu, Aaronson, and Heule constructed a mixed binary--ternary string
rewriting system whose termination is equivalent to the Collatz
conjecture~\cite[Theorem~3.17]{yah}. Their construction turns a question
about integer trajectories into a termination problem for eleven rules.
The system, its arithmetic interpretation, and that equivalence are
their results. This paper studies a restriction on possible
matrix-interpretation proofs for their system.

Rule removal builds a termination proof in stages. An interpretation
orients all remaining rules weakly and a nonempty subset strictly;
relative termination then permits removal of that subset. Repeating
the process can establish termination~\cite[Theorems~2.6 and~2.15]{yah}.
The first step therefore needs at least one strict rule while retaining
weak comparisons for the complete system.

Our result rules out that first step in two coordinates for the
strictly monotone affine format defined below. It quantifies over all
admissible nonnegative real coefficients, including irrational ones,
and imposes no upper bound. A failed bounded SAT search cannot establish
this conclusion: it leaves larger coefficients and coefficients outside
its encoding unexamined. The obstruction is algebraic. It forces all
eleven candidate strict gaps to vanish, rather than excluding one
chosen rule or one coefficient range.

Dimension restrictions matter in this setting. Yolcu, Aaronson, and
Heule report that, for some weakened variants, their tool found only
interpretations of dimension five~\cite[p.~4]{yah}. Their Table~1 also
records a natural dimension-five certificate for the subsystem omitting
\(bf\to ga\) in our notation. These are search outcomes, not proofs
of minimal dimension. The present result supplies a coefficient-independent
exclusion for the full system in dimension two, within a narrower,
precisely stated interpretation format.

The proof has two parts. First, the matrix comparisons force the five
digit matrices to share a triangular orientation, with unit first
diagonals. Second, an affine analysis forces every first-offset gap to
zero. The upper triangular case requires retaining products between
matrix entries and offsets, including several zero-slope cases.
All new obstruction results are formalized in Lean. Appendix~\ref{app:lean-statements}
reproduces the definitions and the exact final theorem type.

\section{The system and the interpretation format}
\label{sec:format}
We use the alphabet $\{a,b,c,d,e,f,g\}$ and the ASCII names from
the authors' implementation~\cite{yahcode}. Table~\ref{tab:symbols}
relates these names to the symbols in their article. The rules are
their full mixed-base system~\cite[Section~3.2]{yah},
with the leftmost symbol interpreted as the outermost map:
\begin{equation}\label{eq:rules}
\begin{array}{lll}
 ad\to d, & bd\to gd, & \\
 ae\to ea, & af\to eb, & ag\to fa,\\
 be\to fb, & bf\to ga, & bg\to gb,\\
 ce\to cb, & cf\to caa, & cg\to cab.
\end{array}
\end{equation}

\begin{table}[ht]
\centering
\begin{tabular}{@{}lccccccc@{}}
\toprule
This paper and the ASCII implementation & $a$ & $b$ & $c$ & $d$ & $e$ & $f$ & $g$\\
Yolcu--Aaronson--Heule~\cite{yah} & $\mathtt f$ & $\mathtt t$ &
$\triangleleft$ & $\triangleright$ & $\mathtt0$ & $\mathtt1$ & $\mathtt2$\\
\bottomrule
\end{tabular}
\caption{Symbol correspondence. The letters $a,b$ are binary digits,
$e,f,g$ are ternary digits, and $c,d$ mark the boundaries.}
\label{tab:symbols}
\end{table}

The five symbols $a,b,e,f,g$ are called digits; $c,d$ are boundary symbols.
Assign to each symbol $s$ an affine map $F_s(x)=M_sx+v_s$ on $\R^2$.
An interpretation is \emph{admissible} when
\begin{equation}\label{eq:admissible}
M_s\in\nn^{2\times2},\qquad v_s\in\nn^2,\qquad (M_s)_{00}\ge1
\quad\text{for every }s.
\end{equation}
For words, composition determines $M_w,v_w$; in particular,
$M_{uv}=M_uM_v$ and $v_{uv}=M_uv_v+v_u$.
A rule $\ell\to r$ is \emph{weak} if $M_\ell\ge M_r$ and $v_\ell\ge v_r$
entrywise. Its first-offset gap is
$\gap(\ell,r)=(v_\ell)_0-(v_r)_0$.
The strictness criterion considered here is $\gap(\ell,r)>0$.

\begin{theorem}\label{thm:main}
For every admissible interpretation in two real coordinates, weak
orientation of all eleven rules in \eqref{eq:rules} implies
$\gap(\ell,r)=0$ for every rule. The same assertion holds for the eleven
rules $\rev(\ell)\to\rev(r)$, where each word is reversed and the direction
of every rule is retained.
\end{theorem}

The two orientations are separate implications, not simultaneous
hypotheses. There is no coefficient bound, integrality requirement,
nonsingularity assumption, or prescribed boundary shape. The conclusion
concerns the first offsets; it does not assert equality of the complete
affine maps. The following sections describe the checked proof. The
formal artifact contains the full algebraic case analyses behind the
compressed arguments below.

\subsection{Scope and soundness}
The theorem covers the strictly monotone, full-context interpretation
format specified in \eqref{eq:admissible} only. Every symbol, including
both boundaries, must satisfy $(M_s)_{00}\ge1$. The relaxed \textsc{top}
setting remains unresolved by this theorem. It does not exclude interpretations used
only for relative top termination, where strict monotonicity of every
symbol is not required. In particular, it does not settle the top-removal
approach in Yolcu--Aaronson--Heule's Lemma~3.18~\cite{yah}.
Arctic interpretations, other strict orders, and higher dimensions are
also outside its scope.

Real coefficients require a fixed positive decrement. For $\delta>0$,
put
\[
x\succ_\delta y\quad\Longleftrightarrow\quad
x_0\ge y_0+\delta\ \text{ and }\ x_1\ge y_1,
\qquad x,y\in\nn^2.
\]
Along a decrease, $\lfloor x_0/\delta\rfloor$ strictly decreases in
$\mathbb N$, so $\succ_\delta$ is well-founded. Nonnegative affine maps
preserve coordinatewise weak order. The bound $(M_s)_{00}\ge1$ also
preserves $\succ_\delta$, since
\[
\bigl(F_s(x)-F_s(y)\bigr)_0
 = (M_s)_{00}(x_0-y_0)+(M_s)_{01}(x_1-y_1)\ge\delta.
\]
A weak rule with first-offset gap $\varepsilon>0$ decreases by at least
$\varepsilon$ at every nonnegative input. For a finite nonempty set of
strict rules, the least such gap supplies a common $\delta>0$.
These observations explain the real-coefficient certificate format;
ordinary strict inequality on the reals alone would not suffice.
Theorem~\ref{thm:main} prevents even this first choice of a strict rule.
The preservation, evaluation, and well-foundedness claims above are
formalized separately in \code{FullTwoSoundness.lean}.

\section{Reduction of the digit matrices}
Write $A,B,C,D,E,F,G$ for the seven matrices and set
$H=A+B+E+F+G$. The matrix parts of the six swaps give
\begin{equation}\label{eq:swaps}
AE\ge EA,\quad AF\ge EB,\quad AG\ge FA,\quad
BE\ge FB,\quad BF\ge GA,\quad BG\ge GB.
\end{equation}

\begin{lemma}\label{lem:matrix}
Under the matrix hypotheses of Theorem~\ref{thm:main}, the five digit
matrices are all upper triangular or all lower triangular in the original
coordinates. Moreover, $A_{00}=B_{00}=E_{00}=F_{00}=G_{00}=1$.
This holds in both word orientations.
\end{lemma}

\begin{proof}[Proof outline]
Suppose $H_{01},H_{10}>0$. With $P=A+B$, the sum of the six nonnegative
differences in \eqref{eq:swaps} is $PH-HP$. The checked identities
$\tr(H^k(PH-HP))=0$, for $k=0,1$, force every entry of this commutator to
vanish: $k=0$ treats the diagonal, and $k=1$ treats the off-diagonal
entries. Each of the six differences therefore vanishes separately.
The resulting exact swaps are incompatible with the boundary
inequalities, as follows.

If $\det F=0$, the exact relations yield a common positive left row.
The rank-one identities and the checked two-dimensional stationary-vector
lemma supply the corresponding right actions. Applying $BD\ge GD$ and
$CF\ge CA^2$ gives $\alpha\ge\beta\ge\alpha^2$, with $\alpha\ge1$.
Thus both eigenvalues are one. Positivity forced by $H_{01},H_{10}>0$
contradicts the rank-one trace and the first-diagonal lower bound.

If $\det F\ne0$, determinant identities give either $\det A=\det B=0$
or nonsingularity of all five digits. For singular $A,B$, their traces
coincide. When $A=B$, the common rank-one argument applies. Otherwise,
writing this positive trace as $\alpha$, exact swaps force
\[
AB=\alpha A,\quad BA=\alpha B,\quad
F=\beta I,\quad E=(\beta/\alpha)A,\quad G=(\beta/\alpha)B.
\]
The boundaries give $\alpha^2\le\beta\le\alpha$; hence
$\alpha=\beta=1$. The diagonal bounds and the two product identities
then force all four mixed off-diagonal products to be zero, contradicting
the assumed positivity of both off-diagonals of $H$.

In the nonsingular case $A=B$ forces $E=F=G$. The inequalities
$AD\ge D$, $ED\le AD$, $CE\ge CA^2$, together with $AE=EA$, give an
entrywise contradiction to $H_{01},H_{10}>0$. If $A\ne B$, an exact
classification gives $\mu,t>0$ with
\begin{gather*}
\tr A=\tr B=3\mu,\qquad \det A=\det B=2\mu^2,\\
E=t(2A-\mu I),\quad F=t(A+B-\mu I),\quad G=t(2B-\mu I).
\end{gather*}
Put $U=A-\mu I$, $V=B-\mu I$. These matrices are nonnegative and satisfy
$U^2=UV=\mu U$, $VU=V^2=\mu V$. The boundary inequalities give
$3t\le2$ and, for $u=(CU)_{00}$, $v=(CV)_{00}$,
\[
u+v\le3tu,\qquad 2\mu(u+v)\le3tv,\qquad u+v>0.
\]
Consequently $\mu\le\tfrac12$. The trace and diagonal bounds force
$\mu=\tfrac12$ and $U_{11}=V_{11}=0$, contradicting
$(U+V)^2=2\mu(U+V)$ and positive off-diagonals.

Thus $H$ has a zero off-diagonal entry. Nonnegativity makes that entry
zero in every digit. Expanding the triangular swap and boundary
inequalities then forces the five first diagonals to be one; both zero
and positive secondary-diagonal cases are included in the formal proof.
Transposition converts the reversed matrix rules to the forward ones.
\end{proof}

\section{The lower triangular affine case}
After Lemma~\ref{lem:matrix}, lower triangular digit maps have the form
\begin{equation}\label{eq:lower}
F_s(x,y)=(x+h_s,\ q_sx+z_sy+k_s),\qquad h_s,q_s,z_s,k_s\ge0.
\end{equation}
Their first offsets add under composition. The six weak swap gaps sum
to zero, so all six vanish and
\begin{equation}\label{eq:twooffsets}
h_a=h_b=r,\qquad h_e=h_f=h_g=u,\qquad r,u\ge0.
\end{equation}

\begin{proposition}\label{prop:lower}
For lower digit maps \eqref{eq:lower}, every first-offset gap vanishes in
either orientation, with arbitrary admissible boundary maps.
\end{proposition}

\begin{proof}[Forward orientation]
The dynamic rule $bd\to gd$ gives $r\ge u$. Write the first row of $C$
as $(c_0,c_1)$, with $c_0\ge1$. If $c_1=0$, the $cf$ boundary gives
$u\ge2r$, hence $r=u=0$, and every boundary gap is zero.

Suppose $c_1>0$. The lower rows of the digit matrix products define
scalar affine maps $y\mapsto z_sy+q_s$. The matrix inequalities induce
all the forward scalar constraints, with outer boundary the identity and
inner boundary the constant $D_{10}/D_{00}$. The checked scalar
classification gives precisely two possibilities: every $q_s=0$, or
every $q_s=v>0$ and every $z_s=0$.

If every $q_s=0$, apply the scalar classification to
$y\mapsto z_sy+k_s$. Put
\[
\Delta_e=k_e-k_b,\quad
\Delta_f=k_f-(1+z_a)k_a,\quad
\Delta_g=k_g-k_a-z_ak_b.
\]
The original boundary offset inequalities give
$c_1\Delta_e\ge c_0(r-u)$ and
$c_1\Delta_f,c_1\Delta_g\ge c_0(2r-u)$.
Since $r\ge u$, these supply the remaining scalar weak constraints.
The scalar result makes all three $\Delta$'s zero. Hence $u\ge2r$,
so $r=u=0$ and all eleven first-offset gaps vanish.

In the other case, $q_s=v>0$ and $z_s=0$. The second-offset $ag$
inequality and first-offset $cf$ inequality yield
\[
k_a-k_f\ge v(r-u),\qquad
k_f-k_a\ge vr+(c_0/c_1)(2r-u).
\]
Their sum forces $r=u=0$. The remaining weak inequalities make all five
$k_s$ equal, so the three boundary gaps vanish as well.
\end{proof}

\begin{proof}[Reversed orientation]
The first-offset swap equalities still give \eqref{eq:twooffsets}; now
the inner boundary $fc\to aac$ gives $u\ge2r$. If $D_{01}=0$, the
outer dynamic comparison gives $r\ge u$, and the result follows.
For $D_{01}>0$, the matrix constraints yield a reversed scalar system
for $z_sy+q_s$. Its checked classification gives either all $q_s=0$ or
\[
z_a=1,\quad z_b=z_f=z_g=0,\quad
q_a=0,\quad q_b=q_f=q_g=\beta\ge0.
\]
In the zero case, the scalar reversed theorem applied to $z_sy+k_s$
gives $k_a=0$ and $k_b=k_g$. The outer $db$ inequality then forces
$r=u=0$. In the exceptional case, $gb\to bg$ gives $k_g\ge k_b$;
combining this with $db\to dg$ again gives $r=u=0$ and $k_b=k_g$.
The remaining swaps force $k_a=0$. The two dynamic gaps are therefore
zero, and the other nine gaps already vanish.
\end{proof}

\section{Upper normalization and aggregate cancellation}
\label{sec:upper-normalization}
Upper digit maps with unit first diagonals have the form
\begin{equation}\label{eq:upper}
F_s(x,y)=(x+q_sy+h_s,\ z_sy+k_s),\qquad q_s,z_s,h_s,k_s\ge0.
\end{equation}
In the forward orientation, replace the boundary maps by
\begin{equation}\label{eq:normal}
\widehat C(x,y)=(x+\kappa y,0),\quad
\widehat D(x,y)=(x,\tau),\qquad
\kappa=C_{01}/C_{00},\quad \tau=(v_d)_1.
\end{equation}
The outer boundary matrix and offset comparisons in coordinate zero
are divided by $C_{00}>0$. The dynamic offset differences are unchanged,
and their normalized matrix differences are zero. Thus all weak rules
are preserved, and zero first-offset gaps in the normalized model imply
zero gaps in the original model. This leaves 22 nonnegative parameters.

For a digit word, composition in \eqref{eq:upper} is described exactly by
\begin{equation}\label{eq:composition}
\begin{aligned}
q_{uv}&=q_v+q_uz_v,& z_{uv}&=z_uz_v,\\
h_{uv}&=h_u+h_v+q_uk_v,& k_{uv}&=k_u+z_uk_v.
\end{aligned}
\end{equation}
In particular, the products $q_uk_v$ remain in every offset comparison.

\begin{lemma}\label{lem:aggregate}
Put
\[
\begin{array}{lll}
p=z_a+z_b-2,& Q=q_a+q_b,& R=k_a+k_b,\\
r=z_e+z_f+z_g-3,& P=q_e+q_f+q_g,& K=k_e+k_f+k_g.
\end{array}
\]
The sums of the six swap gaps in the $q,k,h$ coordinates are
\begin{equation}\label{eq:aggregate}
G_q=rQ-pP\ge0,\quad G_k=pK-rR\ge0,\quad G_h=QK-PR\ge0.
\end{equation}
If $p<0$ or $r<0$, and $Q+P,R+K>0$, all six affine swaps are exact.
\end{lemma}
\begin{proof}
Expanding \eqref{eq:composition} gives \eqref{eq:aggregate} and
\[
RG_q+QG_k=pG_h,\qquad KG_q+PG_k=rG_h.
\]
A negative coefficient on the right forces $G_h=0$. Adding the resulting
equalities gives $(R+K)G_q+(Q+P)G_k=0$, hence $G_q=G_k=0$.
Every individual gap in these sums is nonnegative. The secondary-slope
swap gaps also sum to zero, completing the assertion.
\end{proof}

Two checked degenerate lemmas dispose of $Q+P=0$ and $R+K=0$ without
any condition on the secondary slopes. In the second case all $k_s=0$;
the first offsets reduce to \eqref{eq:twooffsets}. If $\tau=0$,
$bd$ and $cf$ force $r=u=0$. If $\tau>0$, the matrix $q$-constraints
form a reversed scalar system, which forces $q_a=0$ and $q_b=q_g$;
the same conclusion follows. The zero-$q$ case follows from the forward
scalar classification of the $k$-constraints.

The remaining proof therefore assumes positive total $q$ and $k$.
For exact swaps, a common form
$q_s=\lambda(1-z_s)$, $k_s=\eta(1-z_s)$, with $\lambda,\eta>0$,
is impossible. Indeed, $H_s=h_s+\lambda k_s$ has binary value $L$
and ternary value $N$. The dynamic rule gives $L\ge N$, while the $cf$
matrix and offset comparisons give $N\ge2L$. Thus $L=N=0$, forcing
every $k_s=0$, a contradiction. These identities also treat boundary
parameter values equal to zero.

\section{Completion of the upper case and word reversal}
\begin{proposition}\label{prop:upper}
Every normalized upper model satisfying all eleven weak forward rules
has zero first-offset gap in every rule.
\end{proposition}
\begin{proof}[Proof outline]
First suppose $z_a,z_b>0$. The exact slope comparisons make
$z_e=z_f=z_g=s$. If $s>0$, they also make $z_a=z_b=t$.
For exact swaps between triangular matrices with diagonals
$(\alpha,\rho)$ and $(\beta,\sigma)$, with $\alpha,\beta,\rho>0$,
the checked resonance calculation
gives either equal binary off-diagonals and equal ternary off-diagonals,
or $\rho=2\alpha$, $\sigma=3\beta$.

If $\min(s,t)<1$, Lemma~\ref{lem:aggregate} makes the swaps exact.
The $q$ resonance $(t,s)=(2,3)$ is unavailable, so
$q_s=\lambda(1-z_s)$ with $\lambda>0$. The $k$ calculation gives the
common form from Section~\ref{sec:upper-normalization}, except at
$(t,s)=(1/2,1/3)$. At that
exceptional pair, the $ce$ and $cf$ matrix gaps force $\kappa=\lambda$;
the projected offset argument still gives $L\ge N\ge2L$.
If $s,t\ge1$, the unit case is checked separately. For $1\le s\le t$
with $t>1$, the boundary inequalities give $P\ge2Q$, whereas the
aggregate inequality gives $P\le\tfrac32Q$, forcing every $q_s=0$.
For $s>t\ge1$, the $bg$ and $bd$ secondary inequalities force
$k_b=k_g=\tau=0$, and the remaining swaps force every $k_s=0$.
If $s=0$, exact swaps yield the same common form directly, even when
the two positive binary slopes differ.

Now suppose at least one binary slope is zero. If the other exceeds
one, the swap inequalities force every $q_s=0$. Otherwise $p<0$, so
Lemma~\ref{lem:aggregate} again makes all swaps exact. When exactly one
binary slope is zero, the common form follows directly. When both are
zero, the same holds unless
\[
(z_a,z_b,z_e,z_f,z_g)=(0,0,0,1,0),\qquad
q_a=q_b=q_e=q_g=q>0,\quad q_f=0.
\]
In this exceptional branch, $k_e=k_a$, $k_f=0$, $k_g=k_b$, and $cf$
gives $\kappa\ge q$. Exact swaps and the $bd,cf$ offset comparisons
force $h_a+\kappa k_a\le0$, then $k_a=k_b=0$, contradicting positive
total $k$. This exhausts the slope cases; the degenerate lemmas give
all eleven zero-gap equalities in every surviving model.
\end{proof}

For reversed upper models, normalize instead
\[
\widehat D(x,y)=(x+\lambda y,0),\qquad
\widehat C(x,y)=(x,\sigma),\qquad
\lambda=D_{01}/D_{00},\quad \sigma=(v_c)_1.
\]
The same entrywise argument preserves all weak rules and transfers
zero gaps back to the original boundaries. For an upper affine map,
write its homogeneous matrix as
\[
\mathcal H_s=\begin{pmatrix}1&q_s&h_s\\0&z_s&k_s\\0&0&1\end{pmatrix}.
\]
If $J$ reverses the three coordinates, $\Phi(\mathcal H)=J\mathcal H^\mathsf TJ$
reverses products. It exchanges $q_s,k_s$ and preserves $z_s,h_s$.
Thus it preserves entrywise inequalities and the first-offset entry,
while converting the normalized reversed rules into the forward rules.
The boundary forms exchange; the symbol labels do not. This operation
is used only after normalization, when every image is again affine.

\begin{proof}[Proof of Theorem~\ref{thm:main}]
Apply Lemma~\ref{lem:matrix}. For lower digits use
Proposition~\ref{prop:lower}. For upper digits, normalize the boundaries
and use Proposition~\ref{prop:upper}; in the reversed orientation first
apply $\Phi$. The checked normalization lemmas transfer each of the
eleven zero gaps back to the original interpretation.
\end{proof}

\section{Formal statement, verification, and scope}
The formalization uses Lean~4 and mathlib~\cite{lean4,mathlib}. Matrices
and offsets are functions on \code{Fin 2} with values in $\R$.
The shared definition \code{Affine.comp X Y} is
\[
(M_XM_Y,\ M_Xv_Y+v_X),
\]
so the formal composition convention is the one used throughout this
paper. \code{Admissible} is exactly \eqref{eq:admissible};
\code{ForwardWeak} and \code{ReversedRealWeak} each contain eleven
coefficientwise comparisons. The two zero-gap records each contain
eleven equalities at offset index zero.

The main declarations in \code{FullTwoCoordinate.lean} are
\begin{center}
\small
\begin{tabular}{@{}ll@{}}
\toprule
Declaration in \code{CollatzResearch.FullTwo} & Conclusion\\
\midrule
\code{forward\_all\_gaps\_zero} & Eleven forward gaps vanish\\
\code{reversed\_all\_gaps\_zero} & Eleven reversed gaps vanish\\
\code{full\_two\_coordinate\_obstruction} & Both implications\\
\bottomrule
\end{tabular}
\end{center}
Each declaration reports only \code{propext}, \code{Classical.choice},
and \code{Quot.sound}. No additional axiom or proof placeholder is used.
The complete dependency closure, including the soundness lemmas below,
contains 51 local modules. These were rebuilt in a fresh directory, and
an additional audit queries the axioms of every public theorem and lemma,
225 declarations in total. The retained report records the exact source inventory,
source and compiled-module hashes, compiler outputs, dependency revisions,
and the final theorem audit. Existing local compiled files are excluded
from that fresh build.

The fixed-gap soundness lemmas use the namespace
\code{CollatzResearch.FullTwoSoundness}:
\begin{center}\small
\begin{tabular}{@{}ll@{}}
\code{admissible\_preserves\_gap} & Preservation under admissible maps\\
\code{weak\_rule\_gives\_gap} & Evaluation of a weak rule with an offset gap\\
\code{gap\_wellFounded} & Well-foundedness for a positive decrement
\end{tabular}
\end{center}
These three declarations are included in the combined fresh build and
axiom audit. They do not formalize the prior Collatz equivalence or the
general rule-removal theorem.

The checked environment is Lean~4.27.0 with mathlib revision
\begin{center}\small\code{a3a10db0e9d66acbebf76c5e6a135066525ac900}.\end{center}
The main source has SHA-256
\begin{center}\scriptsize\code{429785555eecec277248380fe0c2e767c8c7303cf0e2712de29878dda2bd29a0}.\end{center}
From the artifact root, the verification entry point is
\begin{quote}\small\ttfamily
python3 verify.py \textbackslash\newline
\hspace*{1em}--mathlib-root /path/to/pinned/mathlib \textbackslash\newline
\hspace*{1em}--build-root /path/to/temporary/builds \textbackslash\newline
\hspace*{1em}--report /path/to/fresh-check.json
\end{quote}
The retained verification record is
\path{verification/rebuild.json}. The claim-to-declaration map in
\path{paper/lean-claims.md} identifies the supporting formal results.
Appendix~\ref{app:lean-statements} exposes the actual definitions and
theorem type, so the connection between the prose and the formal claim
can be inspected directly.

The trust boundary consists of Lean's kernel and the stated foundational
axioms, together with the definitions shown in the appendix. Automated
tactics construct proof terms that the kernel checks. The build rejects
proof placeholders and additional axioms; it does not use numerical
search results as proof. Kernel acceptance establishes the Lean
proposition. The symbol table, appendix, and claim-to-declaration map address
the separate question of whether that proposition is the one claimed
in the paper. The Collatz equivalence is cited from prior work and is
not asserted to have been re-formalized by this artifact.

\subsection{Availability and reproducibility}
The paper, Lean sources, verification script, retained compiler output,
and claim-to-declaration map are available at
\begin{center}\small
\url{https://github.com/x1xhlol/collatz-two-coordinate-obstruction}.
\end{center}
The repository records the source hashes and pins the dependencies.
Reproduction uses a new output path to preserve the supplied audit
record. The main obstruction has no coefficient cap, integrality
assumption, or nonsingularity assumption. It does retain every
admissibility condition in Section~\ref{sec:format}. In particular,
the result neither proves the Collatz conjecture nor excludes the
relaxed \textsc{top} setting.

\section*{Acknowledgments}
I thank Emre Yolcu, Scott Aaronson, and Marijn J.\ H.\ Heule for making
their rewriting system and implementation available. The formalization
uses Lean and mathlib, and I thank their developers and contributors.

\section*{Use of AI}
OpenAI Codex was used extensively for mathematical exploration,
development of the Lean proofs and verification scripts, and preparation
and revision of this manuscript. The proof claims are supported by the
accompanying Lean artifacts, rather than by the system's assertions.
This disclosure does not imply review or endorsement by the authors
of the underlying rewriting system.

\begin{thebibliography}{9}\small
\bibitem{yah} E. Yolcu, S. Aaronson, and M. J. H. Heule.
\emph{An Automated Approach to the Collatz Conjecture}.
Journal of Automated Reasoning \textbf{67}(2), Article 15, 2023.
\url{https://doi.org/10.1007/s10817-022-09658-8}.
\bibitem{yahcode} E. Yolcu, S. Aaronson, and M. J. H. Heule.
\emph{Rewriting Collatz: implementation and rewriting systems}.
\url{https://github.com/emreyolcu/rewriting-collatz}.
\bibitem{lean4} L. de Moura and S. Ullrich.
\emph{The Lean 4 Theorem Prover and Programming Language}.
In Automated Deduction, CADE 28, pp. 625--635, 2021.
\url{https://lean-lang.org/papers/lean4.pdf}.
\bibitem{mathlib} The mathlib Community.
\emph{The Lean mathematical library}.
In Certified Programs and Proofs, CPP 2020.
\url{https://doi.org/10.1145/3372885.3373824}.
\end{thebibliography}
\clearpage
\appendix
% Exact excerpts extracted from the frozen Lean sources; no proof body is reproduced.
% Preamble requirements: \usepackage{fontspec,fancyvrb}
%                       \newfontfamily\leanappendixfont{DejaVuSansMono.ttf}
% Insert after \appendix in the main document. This snippet adds no page breaks.
% Source SHA256 ReversedBinaryPowerClosure.lean: 6f448bb1475e4214a10d4acdac8835d526a559d166c93acfd910e39bca23f281
% Source SHA256 ReversedRealNormalization.lean: 947bb86b18f46ba0f150437c4a334c79a18b89bbc73b5855eff9381237b23761
% Source SHA256 FullTwoBasic.lean: 38d6e8029e7fff5a948aca33ad661b28ed232ac527c3627882eac397cf00780e
% Source SHA256 CollatzReversedRealTwoCoordinate.lean: b7cc9035e073d624fd270f6fbde9f2149a290dd93f238643a65aefe4fe434b71
% Source SHA256 FullTwoCoordinate.lean: 429785555eecec277248380fe0c2e767c8c7303cf0e2712de29878dda2bd29a0

\section{Exact Lean declarations}
\label{app:lean-statements}
The blocks below reproduce the checked source verbatim, with the original
line numbers. The theorem blocks end at their types; their proof terms are
omitted. Source paths are relative to the repository root.

The affine definitions and \texttt{ReversedRealWeak} are in
\texttt{CollatzCertificate}. The definitions and theorem with diagonal bounds
are in \texttt{CollatzResearch.FullTwo}. Matrix notation is opened in both
namespaces. The order on vectors is pointwise, so \texttt{Weak X Y} means
that every coefficient of \texttt{Y} is at most the corresponding coefficient
of \texttt{X}. Composition places the left argument outside the right argument.
The unrestricted reversed theorem is in
\texttt{CollatzResearch.RealTwoCoordinate}; it uses only
\texttt{Affine.Nonnegative} and the eleven reversed weak comparisons.

\paragraph{Ambient types and the coefficient order.}
\noindent\texttt{\detokenize{formal/ReversedBinaryPowerClosure.lean}}, lines 10--17.
\begin{Verbatim}[formatcom=\leanappendixfont,fontsize=\scriptsize,numbers=left,firstnumber=10,numbersep=5pt,xleftmargin=18pt]
variable {n : Type*} [Fintype n]

abbrev Vec (n : Type*) := n → ℝ
abbrev Mat (n : Type*) := Matrix n n ℝ

def EntrywiseLE (M N : Mat n) : Prop := ∀ i j, M i j ≤ N i j

local infix:50 " ≤ₘ " => EntrywiseLE
\end{Verbatim}

\paragraph{Affine composition and comparison.}
\noindent\texttt{\detokenize{formal/ReversedBinaryPowerClosure.lean}}, lines 255--266.
\begin{Verbatim}[formatcom=\leanappendixfont,fontsize=\scriptsize,numbers=left,firstnumber=255,numbersep=5pt,xleftmargin=18pt]
structure Affine (n : Type*) where
  matrix : Mat n
  offset : Vec n

def Affine.comp (X Y : Affine n) : Affine n :=
  ⟨X.matrix * Y.matrix, X.matrix *ᵥ Y.offset + X.offset⟩

def Affine.Nonnegative (X : Affine n) : Prop :=
  (0 ≤ₘ X.matrix) ∧ 0 ≤ X.offset

def Affine.Weak (X Y : Affine n) : Prop :=
  (Y.matrix ≤ₘ X.matrix) ∧ Y.offset ≤ X.offset
\end{Verbatim}

\paragraph{Two coordinates and admissibility.}
\noindent\texttt{\detokenize{formal/FullTwoBasic.lean}}, lines 9--11.
\begin{Verbatim}[formatcom=\leanappendixfont,fontsize=\scriptsize,numbers=left,firstnumber=9,numbersep=5pt,xleftmargin=18pt]
abbrev Aff2 := Affine (Fin 2)

def Admissible (X : Aff2) : Prop := X.Nonnegative ∧ 1 ≤ X.matrix 0 0
\end{Verbatim}

\paragraph{Forward weak rules.}
\noindent\texttt{\detokenize{formal/FullTwoBasic.lean}}, lines 21--32.
\begin{Verbatim}[formatcom=\leanappendixfont,fontsize=\scriptsize,numbers=left,firstnumber=21,numbersep=5pt,xleftmargin=18pt]
structure ForwardWeak (A B C D E F G : Aff2) : Prop where
  ad : (A.comp D).Weak D
  bd : (B.comp D).Weak (G.comp D)
  ae : (A.comp E).Weak (E.comp A)
  af : (A.comp F).Weak (E.comp B)
  ag : (A.comp G).Weak (F.comp A)
  be : (B.comp E).Weak (F.comp B)
  bf : (B.comp F).Weak (G.comp A)
  bg : (B.comp G).Weak (G.comp B)
  ce : (C.comp E).Weak (C.comp B)
  cf : (C.comp F).Weak (C.comp (A.comp A))
  cg : (C.comp G).Weak (C.comp (A.comp B))
\end{Verbatim}

\paragraph{Forward offset equalities.}
\noindent\texttt{\detokenize{formal/FullTwoBasic.lean}}, lines 34--45.
\begin{Verbatim}[formatcom=\leanappendixfont,fontsize=\scriptsize,numbers=left,firstnumber=34,numbersep=5pt,xleftmargin=18pt]
structure ForwardGapsZero (A B C D E F G : Aff2) : Prop where
  ad : (A.comp D).offset 0 = D.offset 0
  bd : (B.comp D).offset 0 = (G.comp D).offset 0
  ae : (A.comp E).offset 0 = (E.comp A).offset 0
  af : (A.comp F).offset 0 = (E.comp B).offset 0
  ag : (A.comp G).offset 0 = (F.comp A).offset 0
  be : (B.comp E).offset 0 = (F.comp B).offset 0
  bf : (B.comp F).offset 0 = (G.comp A).offset 0
  bg : (B.comp G).offset 0 = (G.comp B).offset 0
  ce : (C.comp E).offset 0 = (C.comp B).offset 0
  cf : (C.comp F).offset 0 = (C.comp (A.comp A)).offset 0
  cg : (C.comp G).offset 0 = (C.comp (A.comp B)).offset 0
\end{Verbatim}

\paragraph{Reversed weak rules.}
Here the index type is finite and has decidable equality; the final
theorem instantiates it with \texttt{Fin 2}.
\noindent\texttt{\detokenize{formal/ReversedRealNormalization.lean}}, lines 83--94.
\begin{Verbatim}[formatcom=\leanappendixfont,fontsize=\scriptsize,numbers=left,firstnumber=83,numbersep=5pt,xleftmargin=18pt]
structure ReversedRealWeak (A B C D E F G : Affine ι) : Prop where
  da : (D.comp A).Weak D
  db : (D.comp B).Weak (D.comp G)
  ea : (E.comp A).Weak (A.comp E)
  fa : (F.comp A).Weak (B.comp E)
  ga : (G.comp A).Weak (A.comp F)
  eb : (E.comp B).Weak (B.comp F)
  fb : (F.comp B).Weak (A.comp G)
  gb : (G.comp B).Weak (B.comp G)
  ec : (E.comp C).Weak (B.comp C)
  fc : (F.comp C).Weak (A.comp (A.comp C))
  gc : (G.comp C).Weak (B.comp (A.comp C))
\end{Verbatim}

\paragraph{Reversed offset equalities.}
\noindent\texttt{\detokenize{formal/FullTwoBasic.lean}}, lines 47--58.
\begin{Verbatim}[formatcom=\leanappendixfont,fontsize=\scriptsize,numbers=left,firstnumber=47,numbersep=5pt,xleftmargin=18pt]
structure ReversedGapsZero (A B C D E F G : Aff2) : Prop where
  da : (D.comp A).offset 0 = D.offset 0
  db : (D.comp B).offset 0 = (D.comp G).offset 0
  ea : (E.comp A).offset 0 = (A.comp E).offset 0
  fa : (F.comp A).offset 0 = (B.comp E).offset 0
  ga : (G.comp A).offset 0 = (A.comp F).offset 0
  eb : (E.comp B).offset 0 = (B.comp F).offset 0
  fb : (F.comp B).offset 0 = (A.comp G).offset 0
  gb : (G.comp B).offset 0 = (B.comp G).offset 0
  ec : (E.comp C).offset 0 = (B.comp C).offset 0
  fc : (F.comp C).offset 0 = (A.comp (A.comp C)).offset 0
  gc : (G.comp C).offset 0 = (B.comp (A.comp C)).offset 0
\end{Verbatim}

\paragraph{The theorem with diagonal bounds.}
\noindent\texttt{\detokenize{formal/FullTwoCoordinate.lean}}, lines 30--34.
\begin{Verbatim}[formatcom=\leanappendixfont,fontsize=\scriptsize,numbers=left,firstnumber=30,numbersep=5pt,xleftmargin=18pt]
theorem full_two_coordinate_obstruction (A B C D E F G : Aff2)
    (hA : Admissible A) (hB : Admissible B) (hC : Admissible C) (hD : Admissible D)
    (hE : Admissible E) (hF : Admissible F) (hG : Admissible G) :
    (ForwardWeak A B C D E F G → ForwardGapsZero A B C D E F G) ∧
    (ReversedRealWeak A B C D E F G → ReversedGapsZero A B C D E F G)
\end{Verbatim}

\paragraph{The unrestricted reversed theorem.}
\noindent\texttt{\detokenize{formal/CollatzReversedRealTwoCoordinate.lean}}, lines 30--35.
\begin{Verbatim}[formatcom=\leanappendixfont,fontsize=\scriptsize,numbers=left,firstnumber=30,numbersep=5pt,xleftmargin=18pt]
theorem reversed_eligible_offsets_equal
    (A B C D E F G : Affine (Fin 2))
    (hA : A.Nonnegative) (hB : B.Nonnegative) (hC : C.Nonnegative)
    (hD : D.Nonnegative) (hE : E.Nonnegative) (hF : F.Nonnegative) (hG : G.Nonnegative)
    (h : ReversedRealWeak A B C D E F G) :
    (D.comp A).offset = D.offset ∧ (D.comp B).offset = (D.comp G).offset
\end{Verbatim}

\end{document}
